\answer{ex:06:1}{$e=(1-e^{-T_a/\tau_e})e^{-d/\tau_e}$. (а)~$0.110$ проти $4.5\cdot10^{-5}$, у $\approx2.5\cdot10^3$ разів. (б)~$\tau_e^*=T_a/\ln(1+T_a/d)=0.59\,\text{с}$.}
\answer{ex:06:2}{Швидкість $\nu_x\nu_y(A_+\tau_+-A_-\tau_-)=0.8A_+$ на секунду, $\approx2900A_+$ за годину; $\tau_-=\tau_+/0.6=33\ms$.}
\answer{ex:06:3}{$\mathbf e_1$ при $\mu^2+s_1^2>s_2^2$; тут $1.01>0.25$, і нейрон вивчить $\mathbf e_1$, хоча розкид уздовж $\mathbf e_2$ у $25$ разів більший: правило знаходить головний вектор матриці кореляцій, а не коваріацій.}
\answer{ex:06:4}{$V=Re^{-(t_c+L-t)/\tau_\gamma}$ між лампочкою і соком, тож $\dot V=V/\tau_\gamma$; сплеск площею $Re^{-L/\tau_\gamma}$: $0.61R$ і $0.78R$.}
\answer{ex:06:5}{Відношення сигнал/шум $1/(N+1)$, спроб $\propto N+1$: $5\cdot10^5$ спроб $\approx1$~міс. і $5\cdot10^{11}$ спроб $\approx8\cdot10^4$~років.}
\answer{ex:06:6}{(а)~$k_2=1/\tau_d$, $k_1=1/\tau_d-1/\tau_\gamma$. (б)~Хибна помилка $-(V/\tau_\gamma)\,\varepsilon/(1+\varepsilon)$, $\varepsilon=\tau_d/\tau_\gamma$, звідки $\tau_d\le0.105\,\text{с}$.}
\answer{ex:06:7}{$0.46$, $0.90$, $0.99$, $0.95$ і $0.70$ при $d=3\,\text{с}$. Точки лягають на одну криву від частки сліду $F=(1-e^{-T_{\text{cue}}/\tau_e})e^{-d/\tau_e}$: навчання при $F\gtrsim0.1$, випадковий вибір при $F<10^{-2}$.}
\answer{ex:06:8}{$\eta^*=1/\bigl(2\sigma^2(N+2)\bigr)$, за $N+2$ спроби; при $m$ тремтячих з $N$ --- за $N(m+2)/m\ge N+2$: розрідженість не допомагає.}
